problem:  
Let \((M^{2n}, g, J)\) be a compact Kähler manifold with strictly **positive scalar curvature**, but assume that \(M\) does **not** have positive Ricci curvature everywhere (so the Ricci form has mixed sign). Let \((N,h)\) be a compact Riemannian manifold whose universal cover has nonpositive sectional curvature.  

It is well known (via Bochner-type arguments) that if \(M\) has **positive Ricci curvature** and \(N\) has **nonpositive sectional curvature**, then every harmonic map \(u:M\to N\) must be constant. The situation for merely **positive scalar curvature** on a Kähler domain is not covered by classical theorems.  

**Question:**  
Is it true that whenever \(M\) is a compact Kähler manifold with strictly positive scalar curvature and \(N\) has nonpositive sectional curvature, every harmonic map \(u:M\to N\) must still be constant?  
If not, under what additional geometric or complex-analytic conditions on \(M\) (e.g. conditions involving holomorphic sectional curvature, existence of certain Kähler potentials, or structure of the canonical bundle) can one guarantee this rigidity?  

Why is it a "good" problem:  
This problem isolates a subtle gap in the established rigidity theory for harmonic maps: the Bochner method requires positivity of the full Ricci tensor, yet it remains unknown whether the weaker assumption of positive scalar curvature suffices, particularly in the presence of a Kähler structure. Resolving this would deepen our understanding of how curvature positivity interacts with harmonic map energy functionals, and it would bridge complex differential geometry (Kähler curvature decompositions) with global analysis (harmonic map theory). The problem is simply stated but conceptually profound—it asks whether scalar curvature positivity in the complex-geometric setting already enforces harmonic rigidity, a question lying at the frontier between Kähler geometry, complex analysis, and geometric PDEs.